% Source: https://artofproblemsolving.com/community/c3h1727169_algebra_problem
% Stable identity: aops-topic-1727169-q01
The squares up to $100$ are $0,1,4,9,16,25,36,49,64,81,100$. Listing the sums $a^2+b^2\le100$ with $0\le a\le b$ by the larger square $b^2$:
\begin{itemize}
\item $b\le3$: $1,2,4,5,8,9,10,13,18$;
\item $b=4$: $16,17,20,25,32$;
\item $b=5$: $25,26,29,34,41,50$;
\item $b=6$: $36,37,40,45,52,61,72$;
\item $b=7$: $49,50,53,58,65,74,85,98$;
\item $b=8$: $64,65,68,73,80,89,100$;
\item $b=9$: $81,82,85,90,97$;
\item $b=10$: $100$.
\end{itemize}
Removing repeats ($25$, $50$, $65$, $85$, $100$ each appear twice), the distinct values are
\[
1,2,4,5,8,9,10,13,16,17,18,20,25,26,29,32,34,36,37,40,41,45,49,50,52,53,58,61,64,65,68,72,73,74,80,81,82,85,89,90,97,98,100,
\]
which is $\boxed{43}$ numbers.

(This agrees with the two-square theorem: $n$ is a sum of two squares exactly when every prime $p\equiv3\pmod4$ divides $n$ to an even power. If $0$ is not allowed, so that both squares are positive, the answer is $35$: the eight perfect squares $1,4,9,16,36,49,64,81$ that are not also sums of two positive squares drop out.)
